\documentclass[11pt]{article}

\usepackage[a4paper,margin=1in]{geometry}
\usepackage{amsmath,amssymb,amsthm,mathtools}
\usepackage{enumitem}
\usepackage[colorlinks=true,linkcolor=blue,citecolor=blue,urlcolor=blue]{hyperref}

\newtheorem{theorem}{Theorem}
\newtheorem{lemma}{Lemma}
\newtheorem{remark}{Remark}

\newcommand{\one}{\mathbf 1}
\newcommand{\Tr}{\operatorname{Tr}}
\newcommand{\dd}{\,\mathrm d}
\newcommand{\E}{\mathbb E}

\title{A lower bound for odd-cycle densities at edge density at least \texorpdfstring{$2/3$}{2/3}}
\author{}
\date{}

\begin{document}
\maketitle

\begin{abstract}
Let $C_m$ be the cycle of length $m$.  We prove, conditionally on the
Schur--convex path comparison stated in Lemma~\ref{lem:path-comparison},
that if $m=2k+1$ is odd and $W$ is a graphon with edge density
$p=t(K_2,W)\ge 2/3$, then
\[
    t(C_{2k+1},W)\ge p^{2k+1}-p(1-p)^{2k}.
\]
\end{abstract}

\section{Preliminaries}

A \emph{graphon} is a symmetric measurable function
$W:[0,1]^2\to[0,1]$.  If $H$ is a finite graph, its homomorphism density
in $W$ is
\[
    t(H,W)
    :=
    \int_{[0,1]^{V(H)}}
    \prod_{ij\in E(H)} W(x_i,x_j)
    \prod_{v\in V(H)} \dd x_v .
\]
In particular,
\[
    t(K_2,W)=\int_{[0,1]^2}W(x,y)\dd x\dd y.
\]

For a bounded symmetric kernel $U$, let $T_U$ denote the corresponding
Hilbert--Schmidt operator on $L^2[0,1]$:
\[
    (T_U f)(x)=\int_0^1 U(x,y)f(y)\dd y.
\]
We write $J$ for the rank-one operator with kernel identically equal to
$1$; thus
\[
    Jf=\langle f,\one\rangle\one.
\]
For $\ell\ge 0$, define the $\ell$-edge path density of $U$ by
\[
    r_\ell:=t(P_\ell,U)=\langle \one,T_U^\ell\one\rangle,
\]
where $P_\ell$ denotes the path with $\ell$ edges.  Thus
\[
    r_0=1,
    \qquad
    r_1=t(K_2,U).
\]

\section{The Schur--convex path input}

The proof uses the following path comparison inequality.  It is the point
at which Schur convexity enters the argument.

\begin{lemma}[Schur--convex path comparison]
\label{lem:path-comparison}
Let $m=2k+1$ be odd.  Let $U:[0,1]^2\to[0,1]$ be a graphon with
edge density
\[
    q:=t(K_2,U)\le \frac13,
\]
and let $r_\ell=t(P_\ell,U)$.  Define
\[
    A_m(U)
    :=
    \sum_{s=1}^{m}(-1)^{m-s}\frac{m}{s}
    \sum_{\substack{a_1,\dots,a_s\ge 0\\ a_1+\cdots+a_s=m-s}}
    \prod_{i=1}^{s} r_{a_i}.
\]
Then
\[
    A_m(U)-r_{m-1}
    \ge
    (1-q)^m-(1-q)q^{m-1}.
\]
\end{lemma}

\begin{remark}
Lemma~\ref{lem:path-comparison} is the specialised form needed here of
the Schur--convex path inequalities of Kim and Lee.  In their notation,
the key ingredient is that complete homogeneous symmetric functions of
even degree are Schur-convex; after expanding the relevant cyclic word
expression, the difference between the two sides above becomes a
nonnegative linear combination, for $q\le 1/3$, of the even-path deviations
\[
    \sum_{F\in E^+_{2d}(P_n)}t(F,U)-\binom{n}{2d}q^{2d},
\]
which are nonnegative by that Schur--convexity argument.
\end{remark}

\section{Main argument}

\begin{theorem}
Let $m=2k+1$ be odd, and let $W$ be a graphon with edge density
\[
    p:=t(K_2,W)\ge \frac23.
\]
Then
\[
    t(C_m,W)
    \ge
    p^m-p(1-p)^{m-1}.
\]
Equivalently,
\[
    t(C_{2k+1},W)
    \ge
    p^{2k+1}-p(1-p)^{2k}.
\]
\end{theorem}

\begin{proof}
Put
\[
    U:=1-W,
    \qquad
    q:=t(K_2,U)=1-p.
\]
Since $p\ge 2/3$, we have $q\le 1/3$.  We shall prove
\[
    t(C_m,1-U)
    \ge
    (1-q)^m-(1-q)q^{m-1}.
\]

Let $T_U$ be the integral operator associated with $U$.  Since
$T_W=J-T_U$, and since $m\ge 3$, the usual trace formula for cycle
densities gives
\[
    t(C_m,W)=\Tr(T_W^m)=\Tr(J-T_U)^m.
\]
We expand $(J-T_U)^m$ as a sum of noncommutative words in $J$ and $T_U$.

First, the unique word containing no copy of $J$ contributes
\[
    (-1)^m\Tr(T_U^m)=-t(C_m,U),
\]
because $m$ is odd.

Now consider a word containing exactly $s\ge 1$ copies of $J$.  Up to cyclic
rotation inside the trace, such a word has the form
\[
    JT_U^{a_1}JT_U^{a_2}\cdots JT_U^{a_s},
    \qquad
    a_1,\dots,a_s\ge 0,
    \qquad
    a_1+\cdots+a_s=m-s.
\]
Because $Jf=\langle f,\one\rangle\one$, we have
\begin{align*}
    \Tr\bigl(JT_U^{a_1}JT_U^{a_2}\cdots JT_U^{a_s}\bigr)
    &=
    \prod_{i=1}^{s}\langle\one,T_U^{a_i}\one\rangle \\
    &=
    \prod_{i=1}^{s}r_{a_i}.
\end{align*}
For a fixed ordered $s$-tuple $(a_1,\dots,a_s)$, the corresponding cyclic
word has $m/s$ representatives in the linear word expansion.  Therefore
all words containing at least one copy of $J$ contribute exactly
\[
    A_m(U)
    =
    \sum_{s=1}^{m}(-1)^{m-s}\frac{m}{s}
    \sum_{\substack{a_1,\dots,a_s\ge 0\\ a_1+\cdots+a_s=m-s}}
    \prod_{i=1}^{s} r_{a_i}.
\]
Consequently,
\[
    t(C_m,W)=A_m(U)-t(C_m,U).
\]

It remains to replace $t(C_m,U)$ by a path density.  Since $0\le U\le 1$,
deleting one edge from the product defining the cycle density can only
increase the integral.  Hence
\begin{align*}
    t(C_m,U)
    &=
    \int_{[0,1]^m}
    \prod_{i=1}^{m}U(x_i,x_{i+1})
    \prod_{i=1}^{m}\dd x_i \\
    &\le
    \int_{[0,1]^m}
    \prod_{i=1}^{m-1}U(x_i,x_{i+1})
    \prod_{i=1}^{m}\dd x_i \\
    &=
    t(P_{m-1},U)
    =
    r_{m-1},
\end{align*}
where $x_{m+1}=x_1$.  Therefore
\[
    t(C_m,W)
    =A_m(U)-t(C_m,U)
    \ge
    A_m(U)-r_{m-1}.
\]
By Lemma~\ref{lem:path-comparison},
\[
    A_m(U)-r_{m-1}
    \ge
    (1-q)^m-(1-q)q^{m-1}.
\]
Thus
\[
    t(C_m,W)
    \ge
    (1-q)^m-(1-q)q^{m-1}.
\]
Finally, since $q=1-p$, this becomes
\[
    t(C_m,W)
    \ge
    p^m-p(1-p)^{m-1}.
\]
Taking $m=2k+1$ gives
\[
    t(C_{2k+1},W)
    \ge
    p^{2k+1}-p(1-p)^{2k},
\]
as claimed.
\end{proof}

\begin{remark}[Sharpness at Tur\'an densities]
At the Tur\'an densities $p=1-1/r$, $r\ge 3$, equality is attained by the
balanced complete $r$-partite graphon.  Indeed, for this graphon the
nonzero eigenvalues of $T_W$ are
\[
    p=1-\frac1r
    \quad\text{and}\quad
    -\frac1r
    \quad\text{with multiplicity }r-1.
\]
Thus, for odd $m$,
\[
    t(C_m,W)
    =p^m+(r-1)\left(-\frac1r\right)^m
    =p^m-p(1-p)^{m-1}.
\]
\end{remark}

\begin{thebibliography}{9}

\bibitem{KimLee}
Jang Soo Kim and Joonkyung Lee,
\emph{Extended commonality of paths and cycles via Schur convexity},
Journal of Combinatorial Theory, Series B \textbf{166} (2024), 109--122.
Preprint: \href{https://arxiv.org/abs/2210.00977}{arXiv:2210.00977}.

\end{thebibliography}

\end{document}
